\documentclass[../main.tex]{subfiles}
\begin{document}

%\section{Кинематика одномерного движения}

\textit{Фрагменты теории и примеры.} В этом параграфе мы говорим о движении вдоль прямой, иными словами, об одномерном движении. Основы кинематики одномерного движения изложены во многих учебниках. Можно порекомендовать, например, учебник Кикоиных~\cite{Kikoin} (гл.\,1 и гл.\,2) или учебник под редакцией Мякишева~\cite{Balashov} (\S\S~1.1--1.7, 1.20--1.23). Кроме того, многие теоретические вопросы весьма удачно разбираются в статьях~\cite{Kik_art_1, Kik_art_2, Average_velocity, bondarov_1} из журнала <<Квант>>. Поэтому ниже мы не говорим подробно о теории, а только уточняем некоторые определения и разбираем ряд примеров.

При движении тела по прямой $a$ его положение в пространстве в момент времени $t$ обычно описывается его координатой $x(t)$, которая откладывается по координатной оси $OX$, совпадающей с прямой $a$. При решении задачи первым делом следует выбрать направление и ноль координатной оси, исходя из соображений удобства. Когда мы говорим <<выбрать>>, то имеем в виду, что необходимо нарисовать схематичный рисунок и на нём указать: направление оси, положение нуля и координаты важных для решения точек.

Назовём \textit{средней скоростью} быстроту изменения координаты тела. Если в момент времени $t_{1}$ координата равна $x(t_{1})$, а в момент времени $t_{2}$ ($t_{2} > t_{1}$) координата равна $x(t_{2})$, то средняя скорость $\langle v_{x} \rangle$ в промежутке времени от $t_{1}$ до $t_{2}$ равна
\[
\langle v_{x} \rangle= \frac{x(t_{2})- x(t_{1})}{t_{2}-t_{1}}.
\]
Треугольные скобки означают среднее, а индекс <<$x$>> указывает на то, что речь идёт о быстроте изменения координаты $x$. Средняя скорость, определяемая соотношением выше, может быть отрицательной, например, если тело движется в сторону, противоположную положительному направлению оси $OX$. В формулировках задач средней скоростью обычно называют модуль величины:
\[
v_{\text{ср.}} = |\langle v_{x} \rangle|=\left| \frac{x(t_{2})- x(t_{1})}{t_{2}-t_{1}} \right|.
\]
\textit{Средняя путевая скорость} определяется по формуле
\[
v_{\text{ср.пут.}} = \frac{L_{12}}{t_{2}-t_{1}},
\]
где $L_{12}$ --- путь, пройденный телом в промежутке времени от $t_{1}$ до $t_{2}$.

\textit{Мгновенной скоростью} $v_{x}$ мы называем среднюю скорость за очень маленький промежуток времени $\Delta t$:
\[
v_{x}(t) = \frac{x(t+\Delta t)- x(t)}{\Delta t}, \quad \Delta t\rightarrow 0.
\]
Далее мы обычно называем мгновенную скорость просто скоростью. Геометрический смысл формулы состоит в том, что скорость в момент времени $t$ --- это тангенс угла наклона касательной, проведённой к графику зависимости координаты от времени $x(t)$ в точке с абсциссой $t$.

\begin{example}
На рис.\:\ref{fig:pr} вы видите график зависимости координаты некоторого тела от времени. Красные линии --- это отрезки касательных к графику, проведённых в разных точках. Касательная, проведённая в т.\:$D$, практически совпадает с графиком.
\begin{figure}[htbp]
\centering
\includegraphics[]{images/1D/pr.eps} % \textbf{Важно: конвертируйте pr.eps в pr.pdf}
\caption{Зависимость координаты от времени}
\label{fig:pr}
\end{figure}
Продолжая касательные до пересечения с осями координат можно приближённо определить скорость. Значения скоростей равны: $v_{x}(0{,}5\,\hbox{с}) \approx -3~\hbox{м}/\hbox{c}$, $v_{x}(1{,}0\,\hbox{с}) = 0$, $v_{x}(2{,}0\,\hbox{с}) \approx 0{,}75~\hbox{м}/\hbox{c}$, $v_{x}(5{,}0\,\hbox{с}) \approx 1{,}0~\hbox{м}/\hbox{c}$.
\end{example}

Часто бывает весьма непросто точно провести касательную к графику, поэтому способ определения скорости по тангенсу угла наклона касательной к графику координаты следует рассматривать как оценочный, позволяющий, например, выделить участки графика, на которых скорость положительная или отрицательная или сравнить скорости в двух точках. 

При решении задач активно используется график зависимости скорости от времени, поскольку площадь трапеции, ограниченной графиком функции $v_{x}(t)$, осью времени, а также прямыми $t = t_{1}$, $t= t_{2}$ численно равна изменению координаты тела в промежутке времени от $t_{1}$ до $t_{2}$.

\begin{example}
Ломаная $ABCD$ на графике, показанном на рис.\:\ref{fig:pr2}, --- это зависимость скорости материальной точки от времени.
\begin{figure}[htbp]
\centering
\includegraphics[]{images/1D/pr2.eps}
\caption{График зависимости скорости от времени}
\label{fig:pr2}
\end{figure}
Изменение координаты точки к моменту времени $1{,}5t_{0}$ равно площади трапеции $0ABC$:
\[
x(1{,}5t_{0}) - x(0) = S_{0ABC} = \frac{1{,}5t_{0}+t_{0}}{2}\cdot v_{0} = \frac{5v_{0}t_{0}}{4}.
\]
Изменение координаты точки к моменту времени $2t_{0}$ равно разности площадей трапеции $0ABC$ и треугольника $CDE$, поскольку при $t > 1{,}5t_{0}$ скорость точки отрицательная:
\[
x(2t_{0}) - x(0) = S_{0ABC} - S_{CDE}=  v_{0}t_{0}.
\]
Путь, пройденный точкой за время $2t_{0}$, отличается от изменения координаты к этому моменту времени. Действительно, при $t > 1{,}5t_{0}$ координата убывает, а путь продолжает возрастать. Путь, пройденный за время $2t_{0}$ равен сумме площадей трапеции $0ABC$ и треугольника $CDE$:
\[
L(2t_{0}) = S_{0ABC} + S_{CDE}=  2v_{0}t_{0}.
\]
\end{example}

Средняя скорость $\langle v_{x}(t) \rangle$ до момента времени $t_{0}$ остаётся постоянной и равной $v_{0}$, далее начинает уменьшаться и в момент времени $2t_{0}$ становится равна $0{,}5v_{0}$. Средняя путевая скорость до момента времени $1{,}5t_{0}$ изменяется также, как обычная средняя скорость, далее в течение некоторого времени уменьшается, а потом начинает расти и в итоге в момент времени $2t_{0}$ становится равна $v_{0}$.

\begin{example}[АРЗ]
Верблюд удаляется от оазиса в виде круга радиусом $R_{0} = 1~\hbox{км}$, двигаясь по прямой, проходящей через его центр. Скорость верблюда уменьшается обратно пропорционально расстоянию до центра. В начале движения скорость верблюда у границы оазиса равна $v_{0} = 2~\hbox{м}/\hbox{с}$. Определите время движения верблюда до точки, расположенной на расстоянии $8~\hbox{км}$ от границы оазиса.

Идея решения этой задачи состоит в том, чтобы изобразить график зависимости обратной скорости верблюда, иначе говоря, величины $\displaystyle{\frac{1}{v}}$ от расстояния до центра оазиса $r$. Легко видеть, что этот график будет представлять собой отрезок прямой.
\begin{figure}[htbp]
\centering
\includegraphics[]{images/1D/pr3.eps}
\caption{График зависимости $1/v$ от $r$}
\label{fig:pr3}
\end{figure}
Поскольку для скорости верблюда справедливо $v(t) = \displaystyle{\frac{\Delta r}{\Delta t}}$, площадь $\Delta S$ узкого столбика ширины $\Delta r$ и высоты $\displaystyle{\frac{1}{v}}$ равна времени, за которое расстояние изменяется на $\Delta r$:
\[
\Delta S = \frac{1}{v}\cdot \Delta r = \frac{\Delta t}{\Delta r}\Delta r = \Delta t.
\]
Отсюда следует, что время движения от границы оазиса до точки на расстояния $8~\hbox{км}$ от границы (или $9~\hbox{км}$ от центра) равно площади закрашенной трапеции на рис.\:\ref{fig:pr3}. Подсчитывая эту площадь, получим ответ:
\[
T= \frac{0{,}5+4{,}5}{2}\cdot 8000~\hbox{с} = 2\cdot 10^{4}~\hbox{с} \approx 5~\hbox{ч}~30~\hbox{мин}.
\]
\end{example}

При вычислениях с использованием определения мгновенной скорости могут быть полезны следующие математические соотношения, справедливые при $\Delta x\rightarrow 0$:
\[
\frac{(x+\Delta x)^{n}- x^{n}}{\Delta x}=nx^{n-1},
\]
где показатель степени $n$ может быть любым числом: положительным, отрицательным, дробным;
\[
\frac{\sin(x+\Delta x)- \sin x}{\Delta x}=\cos x, \quad \frac{\cos(x+\Delta x)- \cos x}{\Delta x}=-\sin x.
\]

Определение \textit{ускорения} похоже на определение мгновенной скорости:
\[
a_{x}(t) = \frac{v_{x}(t+\Delta t)- v_{x}(t)}{\Delta t}, \quad \Delta t\rightarrow 0.
\]
По физическому смыслу ускорение --- это быстрота изменения скорости. Следует особо отметить, что в общем случае и скорость $v_{x}(t)$, и ускорение $a_{x}(t)$ являются функциями времени. Мы исследуем самые простые модели, в которых либо скорость считается постоянной, либо ускорение. 

При движении с постоянным ускорением скорость изменяется линейно (это следует из определения ускорения), так что справедливо соотношение
\[
v_{x}(t) = v_{x}(0) + a_{x}t.
\]
Зависимость координаты тела $x(t)$ от времени при движении с постоянным ускорением $a_{x}$ даётся формулой
\[
x(t) = x(0) + v_{x}(0)t +\frac{a_{x}t^{2}}{2}.
\]
Комбинируя формулы выше, можно получить полезную формулу для изменения координаты
\[
x(t) - x(0) = \frac{ v_{x}^{2}(t)- v_{x}^{2}(0)}{2a_{x}}.
\]
Движение вдоль вертикальной прямой в поле тяжести Земли при некоторых допущениях (как минимум, следует пренебречь силой сопротивления) --- это типичный случай движения с постоянным ускорением $g$.

\subsection*{Основные задачи}
Везде, где требуется, ускорение свободного падения равно $g = 10\:\hbox{м}/\hbox{с}^{2}$.

\begin{enumerate}[label=\arabic*., leftmargin=*, itemsep=0.5\baselineskip]
\item Поезд въезжает на мост со скоростью $v_{0}$. Если он будет на мосту разгоняться с некоторым ускорением, то проедет мост за время $t_{1} = 30~\hbox{с}$, если с таким же ускорением он будет тормозить, то проедет мост за время $t_{2} = 60~\hbox{с}$. Изобразите качественный график зависимости скорости от времени для каждого случая. За какое время $t_{3}$ поезд проедет мост при равномерном движении со скоростью $v_0$?

\item 1) Небольшое тело, движущееся вдоль прямой, в течение некоторого времени удаляется на расстояние $1~\hbox{м}$ от начальной точки и останавливается. Известно, что ускорение тела в любой момент времени по абсолютной величине не больше $1~\hbox{м}/\hbox{с}^{2}$. Чему равно минимально возможное время движения тела? \\
2) Тело, двигаясь с постоянным ускорением по прямой, удаляется от начальной точки на расстояние $1~\hbox{м}$ за $1~\hbox{с}$. Чему может быть равна скорость $w$ тела в середине пути? В ответе укажите границы диапазона: $w_{\text{min}}$ и $w_{\text{max}}$.

\item Специальный электропоезд отходит от станции, разгоняясь с постоянным ускорением. Первый вагон проходит мимо края платформы за время $t_{\text{перв.}}$, а последний за время $t_{\text{посл.}}$; в начальный момент кабина машиниста находится у самого края платформы. \\
1) Средняя скорость электропоезда (к моменту, когда мимо края платформы пройдут все вагоны) в процессе разгона оказалась равна $u$. Чему равна длина одного вагона? Чему равна длина поезда? \\
2) Во сколько раз отличается время $t_{4}$, за которое четвёртый вагон проходит мимо края платформы, от времени $t_{5}$, за которое мимо края платформы проходит пятый вагон.

\item График зависимости скорости материальной точки от времени показан на рис.\:\ref{fig:vt2}, $v_{0}$ и $t_{0}$ --- известные константы.
\begin{figure}[htbp]
\centering
\includegraphics[]{images/1D/vt2.eps}
\caption{График скорости от времени}
\label{fig:vt2}
\end{figure}
1) Докажите, что максимальная средняя скорость движения точки достигается в момент времени $t_{1}$, такой что площади треугольников $0AB$ и $BCD$ равны. Найдите $t_{1}$.\\
2) Используя любую доступную компьютерную программу, постройте график зависимости средней скорости от времени для интервала времени от $0$ до $2t_{0}$. Для удобства построения можете считать параметры $v_{0}$ и $t_{0}$ безразмерными, равными $1$.\\
3) Представьте, что синяя линия на рис.\:\ref{fig:vt2} --- это график зависимости средней скорости (!) точки от времени. Изобразите для этого случая график зависимости скорости $v(t)$ от времени.

\item Упругий мяч, отпущенный с высоты $h$ над полом, и движущийся всё время вдоль одной прямой, при каждом столкновении с полом теряет часть кинетической энергии. В простейшей модели столкновения скорость мяча в конце столкновения считается равной $v^{\prime} = k\cdot v$, где $v$ --- скорость мяча в начале столкновения, $k$ --- численный коэффициент ($k<1$), который иногда называют \emph{коэффициентом восстановления}. Далее везде в этой задаче следует считать, что время столкновения с полом и сила сопротивления воздуха пренебрежимо малы.\\
1) Пусть коэффициент восстановления при любом столкновении мяча с полом равен $k = 0{,}9$. Мяч отпускают с высоты $h = 1\:\hbox{м}$. На какую высоту поднимется мяч после пятого отскока от пола? Через какое время после начала движения мяч столкнётся с полом тридцатый раз?
\begin{figure}[htbp]
\centering
\includegraphics[width=0.6\linewidth]{images/1D/kicks.eps}
\caption{Амплитудно-временная характеристика ударов}
\label{fig:kicks}
\end{figure}
2) На рис.\:\ref{fig:kicks} вы видите фрагмент зависимости громкости звука от времени при ударах футбольного мяча о пол. График получен с помощью микрофона и программы Audacity. Первый пик соответствует первому удару мяча о пол. Определите коэффициент восстановления $k$ в данном случае. Оцените погрешность полученного результата.

\item Определите экспериментально коэффициент восстановления в процессе многократных столкновений мяча с полом. Запишите звук столкновений с помощью, например, бесплатной программы \href{https://www.audacityteam.org}{Audacity}, а затем, анализируя амплитудно-временную характеристику (следует переключиться в режим <<Волноформа (Дб)>> в меню <<Аудиотрек>>), определите время столкновений. Далее вычислите время между столкновениями, рекомендуется получить не менее трёх точек, по которым можно определить коэффициент восстановления. \\
\textit{Совет.} Произведите эксперимент на кафельном полу. Желательно использовать резиновый мяч <<прыгун>>. Начальную высоту рекомендуем установить порядка $1$ м. Следует позаботиться о том, чтобы в процессе движения мяч как можно меньше смещался по горизонтали от начальной точки.

\item Назовём \emph{безразмерной скоростью} отношение $\displaystyle{u = \frac{v(t)}{v_{0}}}$, где $v(t)$ --- скорость тела в момент времени $t$, а $v_{0}$ --- известная константа. Аналогично, отношение $\tau = \displaystyle{\frac{t}{t_{0}}}$ назовём \emph{безразмерным временем}.
\begin{figure}[htbp]
\centering
\includegraphics[]{images/1D/vt1.eps}
\caption{График безразмерной скорости}
\label{fig:vt1}
\end{figure}
На рис.\,\ref{fig:vt1} линией синего цвета изображена зависимость $u(\tau)$ для некоторого тела, движущегося прямолинейно. Линия красного цвета --- касательная к графику функции $u(\tau)$, проведённая в т.\:$\tau = 0$. \\
1) Чему равно ускорение тела при следующих значениях безразмерного времени: $\tau = 0{,}1$; $\tau = 3{,}9$; $\tau = 1{,}0$?\\
2) Какое расстояние пройдёт тело в промежутке изменения безразмерного времени: от $0$ до $0{,}5$; от $3{,}0$ до $4{,}0$?\\
3) Определите с относительной погрешностью не более $5\:\%$ расстояние, которое проходит тело за период времени от $0$ до $4t_{0}$.\\
4) Зависимость $u(\tau)$, представленная на графике, характерна для падающего (в отсутствие ветра) с большой высоты теннисного мяча. Каков при этом физический смысл отношения $\displaystyle{\frac{v_{0}}{t_{0}}}$ и величины $v_{0}$?

\item Капли воды, образующие туман, --- маленькие (порядка $10\:\hbox{мкм}$), поэтому сила вязкого трения, действующая на каплю со стороны воздуха, оказывается сопоставимой с силой тяжести уже при небольшой скорости падения капли порядка 1 см/с.\\
Синяя линия на рис.\:\ref{fig:vx1} --- это график зависимости скорости $v$ от пройденного расстояния $S$ при вертикальном движении капли вниз; шкала слева. Красная линия --- график зависимости обратной скорости $\displaystyle{\frac{1}{v}}$ от расстояния, пройденного каплей; шкала справа. Ускорение капли в нулевой момент времени обусловлено только силой тяжести и равно $g$.
\begin{figure}[htbp]
\centering
\includegraphics[]{images/1D/vx1.eps}
\caption{Графики скорости и обратной скорости капли}
\label{fig:vx1}
\end{figure}
1) Чему равно ускорение капли на расстоянии порядка метра от начальной точки?\\
2) Найдите ускорение капли на расстоянии $0{,}2\:\hbox{м}$ от начальной точки. \\
3) Определите как можно точнее время, за которое расстояние между каплей и начальной точкой увеличивается от $0{,}2\:\hbox{м}$ до $1{,}0\:\hbox{м}$. Оцените точность полученного результата.

\item Если суммарная мощность сил, действующих на тело, остаётся постоянной при его движении по прямой, то справедливо соотношение
\[
a \cdot v = \text{const},
\]
где $a$ --- ускорение тела в некоторый момент времени, а $v$ --- его скорость в этот момент. Материальная точка движется так, что всё время остаётся справедливым равенство выше, при этом в нулевой момент времени: $a = a_{0} = 1\:\hbox{м}/\hbox{c}^{2}$.\\
1) Изобразите график зависимости величины $\displaystyle{\frac{1}{a}}$ от скорости $v$, если начальная скорость равна: $v_{0} = 0{,}5\:\hbox{м}/\hbox{с}$; $v_{0} = 1\:\hbox{м}/\hbox{с}$; $v_{0} = 2\:\hbox{м}/\hbox{с}$. Для $v_{0} = 1\:\hbox{м}/\hbox{с}$ определите время, за которое скорость увеличивается на $2\:\hbox{м}/\hbox{с}$.\\
2) Известно, что к моменту времени $t = 3\:\hbox{с}$ скорость точки увеличилась в два раза. Чему равна в этом случае скорость $v_{0}$?\\
3) При условиях п.\:2) определите скорость точки в произвольный момент времени $t$.

\item Одинаковые шарики, движущиеся вдоль одной прямой, при абсолютно упругом столкновении обмениваются скоростями.
\begin{figure}[htbp]
\centering
\includegraphics[]{images/1D/two_balls.eps}
\caption{Две одинаковые бусинки до и после столкновения}
\label{fig:two_balls}
\end{figure}
1) Три одинаковые бусинки насажены на гладкую, очень длинную спицу. Для каждой бусинки изобразите качественный график зависимости координаты от времени. Начальную координату и скорость каждой бусинки выберите произвольным образом. Определите максимально возможное количество столкновений между бусинками.
\begin{figure}[htbp]
\centering
\includegraphics[]{images/1D/three_balls.eps}
\caption{Три бусинки на спице}
\label{fig:three_balls}
\end{figure}
2) На гладкую спицу насажено $n$ бусинок. Чему равно максимально возможное количество столкновений между бусинками в этом случае.\\
3) Одинаковые частицы, взаимодействующие упруго при столкновениях, образуют две группы, движущиеся вдоль одной прямой навстречу друг другу. Скорости всех частиц одной группы направлены в одну сторону. В каждой группе скорость $(k+1)$-ой частицы больше скорости $k$-ой частицы.
\end{enumerate}




фон\end{document}